\subsection{Dựng và nhận xét miền elip từng bước}
\label{sec:construct_regions}
Trong hai chiều, đặt $\delta_i=\mu_i-\bar x_i$ và
$r=s_{12}/\sqrt{s_{11}s_{22}}$, với $S\succ0$ nên $|r|<1$.
Ta có
\[
 S^{-1}=\frac1{s_{11}s_{22}-s_{12}^2}
 \begin{pmatrix}s_{22}&-s_{12}\\-s_{12}&s_{11}\end{pmatrix}.
\]
Phương trình miền elip viết tường minh thành
\begin{equation}
 \frac n{1-r^2}\left[
 \frac{\delta_1^2}{s_{11}}
 -\frac{2r\delta_1\delta_2}{\sqrt{s_{11}s_{22}}}
 +\frac{\delta_2^2}{s_{22}}\right]\le c_T^2.
 \label{eq:ellipse_correlation}
\end{equation}
Hạng $\delta_1\delta_2$ cho thấy miền phụ thuộc cách \emph{hai tọa độ
cùng lệch}. Nó không thể nói đầy đủ bằng hai bất đẳng thức riêng lẻ.

\buoc{1}{Xác định tâm và mức tin cậy}
Tâm là $\xbar$, không phải $\vmu_0$. Chọn $1-\alpha$, tính $c_T^2$
theo $n,p$. Một giả thuyết cụ thể chỉ là một điểm để so với miền đã dựng.

\buoc{2}{Xác định các hướng và bán trục}
Giải $S\mathbf e_i=\ell_i\mathbf e_i$, chuẩn hóa $\mathbf e_i$ về độ
dài 1. Bán trục theo hướng đó là $a_i=\sqrt{c_T^2\ell_i/n}$.
Hai đầu mỗi trục là $\xbar\pm a_i\mathbf e_i$.
Trong hai chiều có thể tính trị riêng trực tiếp:
\[
 \ell_{\pm}=\frac{s_{11}+s_{22}
 \pm\sqrt{(s_{11}-s_{22})^2+4s_{12}^2}}2.
\]
Trị riêng lớn cho trục dài: hướng ấy có độ bất định lớn hơn.
Khi $s_{12}=0$, trục elip song song trục tọa độ. Nếu phương sai khác
nhau, elip vẫn không phải đường tròn.
\nguon{Chương 5, §5.4, (5-18), (5-19), tr.~221; hình học dựa trên (4-7), tr.~153}

\begin{example}[Trục của elip trong ví dụ $n=40$]
Với \eqref{eq:intervaldata}, $S=\begin{pmatrix}4&3\\3&9\end{pmatrix}$,
hai trị riêng là $(13\pm\sqrt{61})/2$, xấp xỉ 10.4051 và 2.5949.
Các bán trục 95\% dài khoảng 1.3163 và 0.6573.
Hướng trục dài có thể chọn $\mathbf e_+=(0.4242,0.9056)'$,
nghiêng khoảng $64.9^\circ$ so với chiều dương trục $\mu_1$.
Tương quan dương làm hướng dài nghiêng lên; độ nghiêng không nhất thiết
$45^\circ$ vì hai phương sai 4 và 9 khác nhau.
\end{example}

\buoc{3}{Kiểm tra một điểm và viết nhận xét}
Với điểm ứng viên $\vmu_0$, tính
$q=n(\xbar-\vmu_0)'S^{-1}(\xbar-\vmu_0)$.
Nếu $q\le c_T^2$, điểm nằm trong hoặc trên biên; nếu $q>c_T^2$, điểm
nằm ngoài. Nhận xét nên nêu tâm, hướng biến thiên, độ dài trục và vị trí
của điểm, rồi nối với quyết định Hotelling ở cùng mức $\alpha$.

Một câu nhận xét đủ ý là: ``Elip tập trung quanh trung bình mẫu và kéo
dài theo hướng hai biến cùng tăng. Điểm giả thuyết nằm ngoài elip vì lệch
theo hướng ít biến thiên; do đó kiểm định Hotelling bác bỏ ở mức 5\%.''
Không viết ``điểm nằm ngoài nên dữ liệu chắc chắn sai'', và không gán
xác suất 95\% cho một tham số cố định nằm trong miền đã quan sát.

\subsection{Dựng miền chữ nhật và phân biệt hai cách dựng}
\label{sec:rectangles}
Với hai trung bình thành phần, hình chữ nhật từ hình chiếu $T^2$ là
\begin{equation}
 \mathcal R_T=\prod_{i=1}^2
 \left[\bar x_i-\sqrt{c_T^2}\sqrt{s_{ii}/n};
 \bar x_i+\sqrt{c_T^2}\sqrt{s_{ii}/n}\right].
 \label{eq:rectangle_T}
\end{equation}
Nó chứa toàn bộ elip $\mathcal C$. Bảo đảm phủ của hình chữ nhật thành
phần là \emph{ít nhất} $1-\alpha$; kết quả phủ đúng $1-\alpha$ cho
\emph{mọi tổ hợp} không đồng nghĩa chỉ hai khoảng thành phần cũng có
xác suất phủ đúng bằng mức đó.

Nếu chỉ hai trung bình thành phần được định trước, Bonferroni dùng $m=2$:
\begin{equation}
 \mathcal R_B=\prod_{i=1}^2
 \left[\bar x_i-k_B\sqrt{s_{ii}/n};
 \bar x_i+k_B\sqrt{s_{ii}/n}\right],\quad
 k_B=t_{n-1;1-\alpha/4}.
 \label{eq:rectangle_B}
\end{equation}
Tâm vẫn là $\xbar$; hai nửa cạnh lần lượt là
$h_{B1}=k_B\sqrt{s_{11}/n}$ và $h_{B2}=k_B\sqrt{s_{22}/n}$.
Các cạnh song song với trục tọa độ. Để kiểm tra điểm, chỉ cần
$|\delta_1|\le h_{B1}$ \emph{và} $|\delta_2|\le h_{B2}$.
\nguon{Chương 5, §5.4, hình chiếu của khoảng $T^2$, tr.~226;
 Bonferroni (5-28), (5-29), tr.~232; so sánh các miền trong Example 5.6, tr.~233--234}

\begin{example}[Hai miền không luôn chứa nhau]
Vẫn dùng $n=40$, $\xbar=(10,20)'$, $S$ trong \eqref{eq:intervaldata}.
Lần này danh sách Bonferroni \emph{chỉ có hai trung bình}, khác danh sách
$m=3$ ở ví dụ trước. Với $\alpha=0.05$,
\[
 k_B=t_{39;0.9875}\approx2.3313,\qquad
 (h_{B1},h_{B2})\approx(0.7372,1.1058).
\]
Vậy
\[
 \mathcal R_B=[9.2628;10.7372]\times[18.8942;21.1058].
\]
Xét hai điểm:
\begin{itemize}
 \item $P=(10.70,18.95)'$ có độ lệch $(0.70,-1.05)'$,
 nằm trong Bonferroni nhưng $q=19.6>6.6604$, nên nằm ngoài elip.
 \item $R=(10.80,20.60)'$ có độ lệch $(0.80,0.60)'$,
 có $q=6.4<6.6604$, nên nằm trong elip, nhưng $0.80>0.7372$,
 do đó nằm ngoài Bonferroni.
\end{itemize}
Như vậy cả hai miền đều có bảo đảm tin cậy thích hợp, nhưng hình dạng
và cách đánh giá một điểm khác nhau. Không khẳng định hình chữ nhật
Bonferroni luôn nằm trong elip hoặc luôn chứa elip.
\end{example}

\begin{figure}[htbp]
\centering
\begin{tikzpicture}[x=3.4cm,y=2.3cm,>=Stealth]
 \draw[->] (-1.18,0)--(1.2,0) node[right]{$\delta_1$};
 \draw[->] (0,-1.6)--(0,1.6) node[above]{$\delta_2$};
 \draw[dashed] (-.816114,-1.224171) rectangle (.816114,1.224171);
 \draw[densely dotted,thick] (-.737210,-1.105816) rectangle (.737210,1.105816);
 \draw[domain=0:360,samples=120,smooth,variable=\t]
 plot ({.816114*cos(\t)},{.612085*cos(\t)+1.060163*sin(\t)});
 \fill (.70,-1.05) circle (.025) node[above left]{$P$};
 \fill (.80,.60) circle (.025) node[right]{$R$};
 \fill (0,0) circle (.02) node[below left]{$\xbar$};
\end{tikzpicture}
\caption{Cùng dữ liệu và mức 95\%: nét liền là elip Hotelling,
 nét đứt là hình chữ nhật từ các khoảng thành phần $T^2$, nét chấm là
 Bonferroni cho $m=2$. Trục thể hiện độ lệch khỏi $\xbar$.}
\label{fig:three_regions}
\end{figure}

\begin{center}
\small
\begin{tabular}{L{.21\linewidth}L{.22\linewidth}L{.23\linewidth}L{.23\linewidth}}
\toprule
Tiêu chí & Elip Hotelling & Chữ nhật $T^2$ & Chữ nhật Bonferroni\\
\midrule
Tâm & $\xbar$ & $\xbar$ & $\xbar$\\[3pt]
Thông tin dùng để dựng & Toàn bộ $S$, kể cả $s_{12}$ & $s_{11},s_{22}$ và hệ số $T^2$ & $s_{11},s_{22}$ và hệ số theo $m$\\[3pt]
Hướng biên & Theo vectơ riêng & Song song trục tọa độ & Song song trục tọa độ\\[3pt]
Phủ trong mẫu chuẩn & Đúng $1-\alpha$ cho vectơ & Ít nhất $1-\alpha$ cho các thành phần & Ít nhất $1-\alpha$ cho danh sách định trước\\[3pt]
Nhận xét điểm & Dùng dạng toàn phương & Kiểm tra từng tọa độ & Kiểm tra từng tọa độ\\[3pt]
Quan hệ hình học & Luôn nằm trong $\mathcal R_T$ & Chứa elip & Không có quan hệ chứa nhau cố định với elip\\
\bottomrule
\end{tabular}
\end{center}
Điểm ngoài elip dẫn tới bác bỏ Hotelling. Điểm ngoài Bonferroni vi phạm
ít nhất một khoảng thành phần đã hiệu chỉnh; đó không phải cùng một
quy tắc kiểm định Hotelling, nên các kết luận có thể khác nhau.

\subsection{Vì sao chữ nhật Bonferroni không trực tiếp phản ánh tương quan?}
\label{sec:bonf_correlation}
Công thức \eqref{eq:rectangle_B} không chứa $s_{12}$: mỗi cạnh được dựng
từ một khoảng một chiều cho $\mu_i$. Sau đó Bonferroni kết hợp các
phát biểu bằng một bất đẳng thức đúng với \emph{mọi kiểu phụ thuộc}.
Phương pháp bảo đảm sai số mà không cần ước lượng xác suất chung chính
xác theo tương quan. Vì vậy hình chữ nhật không xoay theo đám mây dữ liệu.

Để thấy rõ, giữ $n=40$, $\xbar=(10,20)'$, mức 95\%, xét
\begin{equation}
 S_\rho=\begin{pmatrix}4&6\rho\\6\rho&9\end{pmatrix},\quad
 \rho=-0.5,\ 0,\ 0.5.
 \label{eq:covariance_family}
\end{equation}
Đây là ba ma trận hiệp phương sai mẫu giả định, cùng phương sai thành
phần nhưng tương quan khác nhau. Trong cả ba, hình chữ nhật Bonferroni
có cùng tâm, cùng nửa cạnh $(0.7372,1.1058)$.
Elip thì đổi hướng và bán trục; tại $\rho=0$, hai bán trục là 0.8161
và 1.2242, còn tại $\rho=\pm0.5$ là 0.6573 và 1.3163.

\begin{figure}[htbp]
\centering
\begin{tikzpicture}[>=Stealth,font=\small]
 \foreach \corr/\shift in {-0.5/0,0/4.6,0.5/9.2}{
  \begin{scope}[xshift=\shift cm,x=1.5cm,y=1.1cm]
   \draw[->] (-1.1,0)--(1.12,0) node[right]{$\delta_1$};
   \draw[->] (0,-1.5)--(0,1.52) node[above]{$\delta_2$};
   \draw[densely dotted,thick] (-.737210,-1.105816) rectangle (.737210,1.105816);
   \draw[domain=0:360,samples=100,smooth,variable=\t]
   plot ({.816114*cos(\t)},
    {1.224171*(\corr*cos(\t)+sqrt(1-\corr*\corr)*sin(\t))});
   \node at (0,-1.9) {$\rho=\corr$};
  \end{scope}
 }
\end{tikzpicture}
\caption{Giữ phương sai thành phần, tâm và thang trục giống nhau:
 nét chấm Bonferroni không đổi; nét liền elip thay đổi theo tương quan.
 Các trục là độ lệch khỏi trung bình mẫu.}
\label{fig:rho_regions}
\end{figure}

\hieu{Hình chữ nhật trả lời ``mỗi trung bình có thể ở khoảng nào?''.
Elip còn trả lời ``hai trung bình có thể ở những vị trí ấy \emph{cùng
nhau} không?''. Với tương quan dương, cặp lệch cùng chiều được đánh giá
khác cặp lệch ngược chiều, dù từng độ lệch riêng có cùng độ lớn.}

Tuy nhiên, không suy rằng ``Bonferroni hoàn toàn không liên quan tương
quan''. Cần phân biệt ba điều:
\begin{enumerate}
 \item \textbf{Công thức biên của hình chữ nhật thành phần} không dùng
 hiệp phương sai chéo. Hình chữ nhật từ các khoảng thành phần $T^2$
 cũng có tính chất này nếu giữ đường chéo, $n,p,\alpha$ cố định.
 \item \textbf{Xác suất phủ đồng thời thực tế} có thể phụ thuộc tương
 quan vì các thống kê một chiều phụ thuộc nhau. Bonferroni chỉ bảo đảm
 một cận dưới, không khẳng định xác suất phủ bằng nhau trong mọi mô hình.
 \item \textbf{Khoảng cho tổ hợp} vẫn dùng $\mathbf a'S\mathbf a$.
 Với hiệu, $\Var$ mẫu là $s_{11}+s_{22}-2s_{12}$; tương quan dương
 làm khoảng hiệu hẹp hơn khi giữ các phương sai thành phần cố định.
\end{enumerate}

\paragraph{Nếu thêm một hiệu thì miền không còn đơn thuần chữ nhật.}
Khi danh sách có $\mu_1,\mu_2,\mu_1-\mu_2$, phải dùng $m=3$ cho
toàn bộ danh sách. Hai khoảng thành phần tạo hình chữ nhật rộng hơn
trường hợp $m=2$. Khoảng hiệu thêm điều kiện
\[
 |(\mu_1-\mu_2)-(\bar x_1-\bar x_2)|
 \le t_{n-1;1-\alpha/6}\sqrt{(s_{11}+s_{22}-2s_{12})/n}.
\]
Đó là một dải chéo; giao của dải với hình chữ nhật thường là đa giác.
Trong ví dụ hiện tại, điều kiện dải là $|\delta_1-\delta_2|\le1.0465$.
Không vẽ riêng hình chữ nhật rồi nói nó thể hiện đầy đủ cả ba khoảng.

\kiemtra{Giữ $s_{11},s_{22}$ và $n$ nhưng đổi dấu $s_{12}$:
 hình chữ nhật thành phần Bonferroni có đổi không? Không.
 Elip có đổi không? Có, hướng nghiêng đổi.
 Khoảng cho $\mu_1-\mu_2$ có đổi không? Có, vì có hạng $-2s_{12}$.}
